\Article{4}
Bezeichnen wir die relative Wahrscheinlichkeit des Gesamtfehlers \(x\) bei einer bestimmten Beobachtungsart durch die Charakteristik \(\varphi x\), so ist dies wegen der Stetigkeit der Fehler so zu verstehen, dass die Wahrscheinlichkeit eines Fehlers zwischen den unendlich nahen Grenzen \(x\) und \(x+\mathrm{d}x\) gleich \(=\varphi x.\mathrm{d}x\) ist. In der Praxis wird es kaum, ja kaum jemals überhaupt möglich sein, diese Funktion a priori anzugeben: Dennoch lassen sich mehrere allgemeine Ergebnisse über sie feststellen, die wir im Folgenden vorlegen werden. Folgendes ist offensichtlich. Die Funktion \(\varphi x\) ist insofern zu den unstetigen Funktionen zu rechnen, als sie für alle Werte von \(x\) außerhalb der Grenzen der möglichen Fehler gleich \(=0\) sein muss; innerhalb dieser Grenzen wird sie dagegen überall einen positiven Wert annehmen (wenn wir den am Ende des vorigen Artikels besprochenen Fall außer Betracht lassen). In den meisten Fällen wird man positive und negative Fehler derselben Größe als gleich wahrscheinlich voraussetzen dürfen, sodass \(\varphi(-x)=\varphi x\) sein wird. Da ferner kleinere Fehler leichter begangen werden als größere, wird der Wert von \(\varphi x\) meistens für \(x=0\) am größten sein und stetig abnehmen, während \(x\) wächst.

Im Allgemeinen wird aber der Wert des Integrals \(\int\varphi x.\mathrm{d}x\), von \(x=a\) bis \(x=b\) erstreckt, die Wahrscheinlichkeit ausdrücken, dass ein noch unbekannter Fehler zwischen den Grenzen \(a\) und \(b\) liegt. Der Wert dieses Integrals von der unteren Grenze aller möglichen Fehler bis zur oberen Grenze wird daher stets \(=1\) sein. Und da \(\varphi x\)
\SourceLeaf{011}für alle Werte von \(x\) außerhalb dieser Grenzen stets \(=0\) ist, wird offensichtlich auch

\textit{der Wert des Integrals \(\int\varphi x.\mathrm{d}x\) von \(x=-\infty\) bis \(x=+\infty\) erstreckt stets \(=1\) sein.}

\Article{5}
Betrachten wir ferner das Integral \(\int x\varphi x.\mathrm{d}x\) zwischen denselben Grenzen, dessen Wert wir \(=k\) setzen wollen. Sind alle einfachen Fehlerursachen so beschaffen, dass kein Grund vorliegt, weshalb von zwei gleich großen, aber mit entgegengesetzten Vorzeichen versehenen Fehlern der eine leichter entstehen sollte als der andere, so wird dies auch für den Gesamtfehler gelten, oder es wird \(\varphi(-x)=\varphi x\) und folglich notwendig \(k=0\) sein. Hieraus schließen wir, dass, wenn \(k\) nicht verschwindet, sondern z. B. eine positive Größe ist, notwendig die eine oder andere Fehlerursache vorhanden sein muss, die entweder nur positive Fehler erzeugen kann oder wenigstens positive Fehler leichter als negative. Diese Größe \(k\), die tatsächlich das Mittel aller möglichen Fehler oder der Mittelwert von \(x\) ist, kann zweckmäßig der konstante Teil des Fehlers genannt werden. Im Übrigen lässt sich leicht beweisen, dass der konstante Teil des Gesamtfehlers gleich der Summe der konstanten Teile ist, welche die aus den einzelnen einfachen Ursachen hervorgehenden Fehler enthalten. Wird die Größe \(k\) als bekannt vorausgesetzt und von jeder Beobachtung abgezogen und der Fehler der so berichtigten Beobachtung mit \(x'\), seine Wahrscheinlichkeit aber mit \(\varphi' x'\) bezeichnet, so wird \(x'=x-k\), \(\varphi' x'=\varphi x\) und folglich \(\int x'\varphi' x'.\mathrm{d}x'=\int x\varphi x.\mathrm{d}x-\int k\varphi x.\mathrm{d}x=k-k=0\) sein, d. h. die Fehler der berichtigten Beobachtungen werden keinen konstanten Teil haben, was auch an sich offensichtlich ist.

\Article{6}
Wie das Integral \(\int x\varphi x.\mathrm{d}x\) oder der Mittelwert von \(x\) entweder das Fehlen eines konstanten Fehlers oder dessen Vorhandensein und Größe erkennen lässt, so scheint das Integral
\[
\int xx\varphi x.\mathrm{d}x
\]
von \(x=-\infty\) bis \(x=+\infty\) erstreckt (oder der Mittelwert des Quadrats \(xx\)) am geeignetsten, die Unsicherheit der Beobachtungen im Allgemeinen zu bestimmen und zu messen, sodass von zwei Beobachtungssystemen, die sich hinsichtlich der Wahrscheinlichkeit ihrer Fehler unterscheiden, diejenigen als genauer gelten sollen, bei denen das Integral \(\int xx\varphi x.\mathrm{d}x\) den kleineren Wert erhält. Sollte jemand einwenden, dass diese Art nach Willkür, ohne zwingende Notwendigkeit gewählt worden sei, so werden wir bereitwillig zustimmen.
\SourceLeaf{012}Denn diese Frage enthält der Natur der Sache nach etwas Unbestimmtes, das nur durch einen einigermaßen willkürlichen Grundsatz in Grenzen gefasst werden kann. Die Bestimmung einer Größe durch eine Beobachtung, die einem größeren oder kleineren Fehler ausgesetzt ist, lässt sich nicht unpassend mit einem Spiel vergleichen, bei dem es nur Verluste und keine Gewinne gibt, da jeder zu befürchtende Fehler einem Verlust ähnlich ist. Die Kosten eines solchen Spiels werden nach dem wahrscheinlichen Verlust bemessen, nämlich nach der Summe der Produkte aus den einzelnen möglichen Verlusten und ihren jeweiligen Wahrscheinlichkeiten. Welchem Verlust aber jeder Beobachtungsfehler gleichzusetzen ist, ist keineswegs an sich offensichtlich; vielmehr hängt diese Bestimmung zum Teil von unserer Willkür ab. Den Verlust dem Fehler selbst gleichzusetzen, ist offensichtlich unzulässig; denn würden die positiven Fehler als Verluste genommen, so müssten die negativen Gewinne darstellen. Die Größe des Verlustes muss vielmehr durch eine solche Funktion des Fehlers ausgedrückt werden, die ihrer Natur nach stets positiv ist. Da es unendlich viele derartige Funktionen gibt, scheint die einfachste mit dieser Eigenschaft allen anderen vorzuziehen zu sein, und diese ist unstreitig das Quadrat: So ergibt sich der oben vorgelegte Grundsatz.

Der berühmte \textsc{Laplace} hat die Sache zwar auf ähnliche Weise betrachtet, aber den Fehler selbst, stets positiv genommen, als Maß des Verlustes angenommen. Doch wenn wir uns nicht irren, ist diese Art wenigstens nicht weniger willkürlich als unsere: Ob nämlich ein doppelter Fehler als ebenso erträglich wie ein einfacher, zweimal wiederholter oder als schwerer erträglich angesehen wird und ob es daher angemessener ist, dem doppelten Fehler nur das doppelte oder ein größeres Gewicht beizulegen, ist eine Frage, die weder an sich offensichtlich noch durch mathematische Beweise zu entscheiden ist, sondern allein der freien Wahl überlassen werden muss. Außerdem lässt sich nicht leugnen, dass bei jener Art die Stetigkeit verletzt wird: Und gerade deshalb widerstrebt jene Methode der analytischen Behandlung stärker, während sich die Ergebnisse, zu denen unser Grundsatz führt, durch ihre bewundernswerte Einfachheit und Allgemeinheit empfehlen.

\clearpage
